\Article{16}
In artt. 12---15 we supposed that the errors \(x,x',x''\) etc. belonged to the same kind of observations, so that their individual probabilities were expre\SourceLeaf{021}ssed by the same function. But it is evident of itself that the general investigation of artt. 12---14 extends just as easily to the more general case in which the probabilities of the errors \(x,x',x''\) etc. are expressed by different functions \(\varphi x,\varphi'x',\varphi''x''\) etc., i.e. in which those errors belong to observations of different precision or uncertainty. Let us suppose that \(x\) is the error of an observation whose mean error to be feared is \(=m\); and that \(x',x''\) etc. are the errors of other observations whose mean errors to be feared are respectively \(m',m''\) etc. Then the mean value of the sum \(xx+x'x'+x''x''+\text{etc.}\) will be \(mm+m'm'+m''m''+\text{etc.}\). Now if it is known from elsewhere that the quantities \(m,m',m''\) etc. are in a given ratio, say respectively proportional to the numbers \(1,\mu',\mu''\) etc., the mean value of the expression
\[
\frac{xx+x'x'+x''x''+\text{etc.}}{1+\mu'\mu'+\mu''\mu''+\text{etc.}}
\]
will be \(=mm\). But if we put a definite value of the same expression, according as chance supplies the errors \(x,x',x''\) etc., equal to \(mm\), the mean error to which this determination remains subject is found in a similar way to that in the preceding article
\[
=\frac{\BareRoot{(n^4+n'{}^4+n''{}^4+\text{etc.}-m^4-m'{}^4-m''{}^4-\text{etc.})}}{1+\mu'\mu'+\mu''\mu''+\text{etc.}}
\]
where \(n',n''\) etc. are supposed to signify, with regard to the observations to which the errors \(x',x''\) etc. belong, the same thing as \(n\) with regard to the first observation. If, therefore, it is permissible to suppose the numbers \(n,n',n''\) etc. proportional to \(m,m',m''\) etc., that mean error to be feared becomes
\[
=\frac{\BareRoot{(n^4-m^4)}.\BareRoot{(1+\mu'{}^4+\mu''{}^4+\text{etc.})}}{1+\mu'\mu'+\mu''\mu''+\text{etc.}}
\]

But this way of determining an approximate value of \(m\) is not the most suitable. To show this more clearly, let us consider the more general expression
\[
y=\frac{xx+\alpha'x'x'+\alpha''x''x''+\text{etc.}}{1+\alpha'\mu'\mu'+\alpha''\mu''\mu''+\text{etc.}}
\]
whose mean value will likewise be \(=mm\), however the coefficients \(\alpha',\alpha''\) etc. are chosen. The mean error to be feared when we suppose a definite value of \(y\), according as chance supplies the errors \(x,x',x''\) etc., equal to \(mm\) is found by the principles given above
\[
=\frac{\BareRoot{(n^4-m^4+\alpha'\alpha'(n'{}^4-m'{}^4)+\alpha''\alpha''(n''{}^4-m''{}^4)+\text{etc.})}}{1+\alpha'\mu'\mu'+\alpha''\mu''\mu''+\text{etc.}}
\]
\SourceLeaf{022}To make this mean error as small as possible, we shall have to put
\[
\begin{aligned}
\alpha'&=\frac{n^4-m^4}{n'{}^4-m'{}^4}.\mu'\mu'\\[4pt]
\alpha''&=\frac{n^4-m^4}{n''{}^4-m''{}^4}.\mu''\mu''\text{ etc.}
\end{aligned}
\]
Evidently these values cannot be worked out unless the ratio of the quantities \(n,n',n''\) etc. to \(m,m',m''\) etc. is also known from elsewhere: in the absence of this exact knowledge, it at least seems safest\footnote{Namely, we conceive that knowledge of the quantities \(\mu',\mu''\) etc. is within our power only in the case where, by the nature of the matter, the errors \(x,x',x''\) etc. proportional to \(1,\mu',\mu''\) etc. must be regarded as equally probable, or rather where
\[
\varphi x=\mu'\varphi'(\mu'x)=\mu''\varphi''(\mu''x)\text{ etc.}
\]
} to suppose the former proportional to the latter (see art. 11), whence the values result
\[
\alpha'=\frac1{\mu'\mu'},\quad\alpha''=\frac1{\mu''\mu''}\text{ etc.}
\]
i.e. the coefficients \(\alpha',\alpha''\) etc. must be put equal to the relative weights of the observations to which the errors \(x',x''\) etc. belong, taking the weight of the observation to which the error \(x\) belongs as the unit. In this way, with \(\sigma\) denoting as above the number of the given errors, we shall have the mean value of the expression
\[
\frac{xx+\alpha'x'x'+\alpha''x''x''+\text{etc.}}{\sigma}
\]
\(=mm\), and the mean error to be feared when we adopt a value of this expression determined by chance as the true value of \(mm\)
\[
\frac{\BareRoot{(n^4+\alpha'\alpha'n'{}^4+\alpha''\alpha''n''{}^4+\text{etc.}-\sigma m^4)}}{\sigma}
\]
and hence, if indeed it is permissible to suppose \(n,n',n''\) etc. proportional to \(m,m',m''\),
\[
\mathord{\cdot}=\BareRoot{\frac{n^4-m^4}{\sigma}}
\]
which formula is identical with the one we found above for the case of observations of the same kind.

\Article{17}
If the value of a quantity that depends on another unknown quantity has been determined by an obser\SourceLeaf{023}vation that does not have absolute precision, the value of the unknown calculated from it will also be subject to error, but nothing in this determination is left to arbitrary choice. But if \textit{several} quantities depending on the same unknown have become known through observations that are not absolutely exact, we can obtain the value of the unknown either from any one of these observations or from some combination of several observations, which can be done in infinitely many different ways. Although the value of the unknown arising in this way always remains subject to error, a greater error will be feared in one combination, a smaller in another. The case will be similar if several quantities depending simultaneously on several unknowns have been observed: according as the number of observations is equal to, less than or greater than the number of unknowns, the problem will be determined, undetermined or more than determined (at least generally speaking), and in the third case the observations can be combined in infinitely many different ways to determine the unknowns. To select from such a variety of combinations those most suited to the purpose, i.e. those giving values of the unknowns subject to the smallest errors, is certainly by far the most important problem in the application of mathematics to natural philosophy.

In Theoria motus corporum coelestium we showed how the \textit{most probable} values of the unknowns are to be obtained if the law of probability of the errors of observations is known; and since this law by its nature remains hypothetical in almost all cases, we applied that theory to the most plausible law, where the probability of the error \(x\) is supposed proportional to the exponential quantity \(e^{-hhxx}\). From this arose the method long used by us in calculations, especially astronomical ones, and now indeed used by most people doing calculations under the name of the method of least squares.

Later the illustrious \textsc{Laplace}, approaching the matter in another way, showed that the same principle must still be preferred to all others, whatever the law of probability of the errors, provided only that the number of observations is very large. But for a moderate number of observations the matter remained untouched, so that, if our hypothetical law is rejected, the method of least squares must be considered preferable to others only on the ground that it is particularly suited to convenient calculation.

We therefore hope that geometers will welcome our showing, in this new treatment of the subject, that the method of least squares gives the best combination of all. Not approximately, but absolutely, whatever the law of probability of the errors, whatever the number of observations, provided only that we establish the notion of \SourceLeaf{024}mean error according to our treatment in artt. 5 and 6, rather than according to the view of the illustrious \textsc{Laplace}.

Moreover, it is appropriate to warn here in express words that all the following investigations concern only irregular errors free of a constant part, since the complete art of observation properly requires removing all causes of constant errors with the greatest care. What aids the calculus of probabilities itself can offer to someone undertaking the treatment of observations that there is good reason to suspect are not free of constant errors, we reserve for a separate investigation to be published on another occasion.

\clearpage
